\chapter{Sine-Gordon, Massive Thirring, And Affine Toda Theories}
\label{ch:integrable-sine-gordon-thirring-affine-toda}

\ConventionsMixedCFT

This chapter introduces exact relativistic model families in which the
factorized-scattering, form-factor, and TBA structures of the preceding
chapters can be tested against Lagrangian and ultraviolet data.  Every
normalization below is part of the definition of the model being discussed.

\section{Model-Family Datum}

An integrable model family is specified by:
\[
  (\mathcal A_{\mathrm{UV}},\mathcal O_{\mathrm{pert}},
  \lambda,\mathcal Q,\mathcal S,\mathcal F),
\]
where \(\mathcal A_{\mathrm{UV}}\) is a short-distance operator algebra or
regulated Lagrangian chart, \(\mathcal O_{\mathrm{pert}}\) is the relevant
perturbing coordinate, \(\lambda\) is its coupling, \(\mathcal Q\) is the
commutative algebra of conserved charges, \(\mathcal S\) is the factorized
scattering datum, and \(\mathcal F\) is the form-factor datum for local
operators.  The model family is under control only when these entries are
matched in common conventions.

\section{Sine-Gordon Normalization}

Let \(\varphi\) be a real scalar field with Lorentzian action
\[
  S_{\mathrm{SG}}
  =
  \int d^2x\,
  \left[
  -\frac12 \partial_\mu\varphi\,\partial^\mu\varphi
  +\frac{\mu}{\beta^2}\bigl(\cos(\beta\varphi)-1\bigr)
  \right],
\]
where \(0<\beta^2<8\pi\) in the short-distance free-boson normalization used
for the displayed formula.  The potential is periodic under
\[
  \varphi\longmapsto \varphi+\frac{2\pi}{\beta}.
\]
Finite-energy configurations on the line therefore have topological charge
\[
  Q_{\mathrm{top}}
  =
  \frac{\beta}{2\pi}
  \int_{-\infty}^{\infty}dx\,\partial_x\varphi
  \in\mathbb Z .
\]
The charge \(Q_{\mathrm{top}}=1\) particle is the soliton, and
\(Q_{\mathrm{top}}=-1\) is the antisoliton.

\section{Sine-Gordon Spectrum}

Introduce
\[
  \xi=\frac{\beta^2}{8\pi-\beta^2}.
\]
Equivalently \(\lambda=\xi^{-1}=8\pi/\beta^2-1\).  The attractive regime is
\(0<\xi<1\), or \(\lambda>1\).  In this regime the soliton-antisoliton
S-matrix has bound-state poles corresponding to neutral breathers
\[
  B_n,\qquad
  n=1,\ldots,N_B,\qquad
  N_B=\max\{n\in\mathbb Z_{>0}\mid n\xi<1\},
\]
with masses
\[
  m_n=2M_s\sin\!\left(\frac{\pi n\xi}{2}\right),
\]
where \(M_s\) is the soliton mass.  The formula follows from relativistic
two-body kinematics: a pole at rapidity difference \(\theta=\ii u\) in a
two-particle channel of equal masses \(M_s\) produces a bound state of mass
\[
  m^2=2M_s^2(1+\cos u)
  =
  4M_s^2\cos^2\!\left(\frac{u}{2}\right).
\]
For the breather poles \(u=\pi(1-n\xi)\), this gives the displayed value.

\section{Exact Soliton-Antisoliton Scattering Datum}

The one-soliton internal space is
\[
  V=\C\ket{+}\oplus\C\ket{-},
  \qquad
  Q_{\mathrm{top}}\ket{\pm}=\pm\ket{\pm}.
\]
Charge conservation forces the two-particle scattering operator on
\(V\otimes V\) to have the block form
\begin{equation}
  \mathsf S_{\mathrm{sol}}(\theta)
  =
  S_{\mathrm{min}}(\theta)
  \begin{pmatrix}
    1&0&0&0\\
    0&b(\theta)&c(\theta)&0\\
    0&c(\theta)&b(\theta)&0\\
    0&0&0&1
  \end{pmatrix}_{\ket{++},\ket{+-},\ket{-+},\ket{--}},
  \label{eq:sine-gordon-soliton-matrix}
\end{equation}
where
\[
  b(\theta)
  =
  \frac{\sinh(\lambda\theta)}
       {\sinh(\lambda(\ii\pi-\theta))},
  \qquad
  c(\theta)
  =
  \frac{\ii\sin(\pi\lambda)}
       {\sinh(\lambda(\ii\pi-\theta))},
  \qquad
  \lambda=\xi^{-1}.
\]
One standard minimal scalar normalization is the boundary value, for real
\(\theta\), of
\begin{equation}
  S_{\mathrm{min}}(\theta)
  =
  -\exp\left[
    -\ii\int_0^\infty\frac{\dd t}{t}\,
    \frac{\sinh\!\left((1-\xi)t/2\right)}
         {\sinh(\xi t/2)\cosh(t/2)}
    \sin\!\left(\frac{\theta t}{\pi}\right)
  \right],
  \label{eq:sine-gordon-minimal-scalar-factor}
\end{equation}
continued meromorphically from the physical line with the prescribed
physical-strip pole data.  The displayed datum is an on-shell bootstrap
datum: the construction of a local Hilbert-space QFT whose soliton sectors
realize it is a separate existence problem, recorded at the end of the
chapter.

\begin{proposition}[Matrix unitarity and the breather pole locations]
\label{prop:sine-gordon-soliton-matrix-unitarity}
For real \(\lambda>0\), the matrix factor in
\eqref{eq:sine-gordon-soliton-matrix} obeys
\[
  R_\lambda(\theta)R_\lambda(-\theta)=\mathbf 1
\]
as a meromorphic \(4\times4\) matrix.  In the attractive regime
\(\lambda>1\), the denominators of \(b(\theta)\) and \(c(\theta)\) have
physical-strip zeros at
\begin{equation}
  \theta=\ii u_n,\qquad
  u_n=\pi(1-n\xi),
  \qquad n\in\mathbb Z_{>0},\quad n\xi<1,
  \label{eq:sine-gordon-breather-pole-locations}
\end{equation}
and these are precisely the rapidity locations that give the breather masses
\[
  m_n=2M_s\sin\!\left(\frac{\pi n\xi}{2}\right).
\]
\end{proposition}

\begin{proof}
Only the charge-zero block is nontrivial.  Set
\[
  X=\sinh(\lambda\theta),
  \qquad
  s=\sin(\pi\lambda),
  \qquad
  D_\pm=\sinh(\lambda(\ii\pi\mp\theta)).
\]
Using
\[
  D_+D_-
  =
  \sinh(\ii\pi\lambda-\lambda\theta)
  \sinh(\ii\pi\lambda+\lambda\theta)
  =
  -\bigl(X^2+s^2\bigr),
\]
one obtains
\[
  b(\theta)b(-\theta)+c(\theta)c(-\theta)
  =
  \frac{-X^2-s^2}{D_+D_-}=1,
\]
and
\[
  b(\theta)c(-\theta)+c(\theta)b(-\theta)
  =
  \frac{X\,\ii s+\ii s(-X)}{D_+D_-}=0.
\]
The same calculation applies to the second off-diagonal entry, while the
\(\ket{++}\) and \(\ket{--}\) channels are immediate.

The denominator vanishes exactly when
\[
  \lambda(\ii\pi-\theta)=\ii\pi n,\qquad n\in\mathbb Z,
\]
so \(\theta=\ii\pi(1-n/\lambda)=\ii\pi(1-n\xi)\).  This point lies in the
physical strip \(0<\operatorname{Im}\theta<\pi\) exactly when \(n\xi<1\).
Substituting \(u_n=\pi(1-n\xi)\) into the equal-mass bound-state formula gives
\[
  m_n
  =
  2M_s\cos\!\left(\frac{u_n}{2}\right)
  =
  2M_s\sin\!\left(\frac{\pi n\xi}{2}\right),
\]
as claimed.
\end{proof}

\begin{remark}[Free-fermion point]
At \(\beta^2=4\pi\), one has \(\xi=1\) and \(\lambda=1\).  Then
\(\sin(\pi\lambda)=0\), so the reflection amplitude \(c(\theta)\) vanishes,
while \(b(\theta)=1\).  The remaining scalar factor is
\(S_{\mathrm{min}}(\theta)=-1\), the fermionic exchange phase of the massive
Thirring field at \(g_{\mathrm T}=0\).  This is a useful convention check:
the soliton and antisoliton then form the free massive Dirac particle and
antiparticle.
\end{remark}

\section{Lightest Breather Scattering}

The lightest breather \(B_1\) is a neutral scalar particle in the attractive
regime \(0<\xi<1\).  Its diagonal two-body amplitude is
\begin{equation}
  S_{11}(\theta)
  =
  \frac{\sinh\theta+\ii\sin(\pi\xi)}
       {\sinh\theta-\ii\sin(\pi\xi)} .
  \label{eq:sine-gordon-b1b1-smatrix}
\end{equation}
This formula is a useful test case because each analytic requirement can be
checked directly.  Unitarity follows by multiplication:
\[
  S_{11}(\theta)S_{11}(-\theta)
  =
  \frac{\sinh\theta+\ii s_\xi}{\sinh\theta-\ii s_\xi}
  \frac{-\sinh\theta+\ii s_\xi}{-\sinh\theta-\ii s_\xi}
  =
  1,
  \qquad s_\xi=\sin(\pi\xi).
\]
Hermitian analyticity follows from
\[
  S_{11}(\theta)^\ast=S_{11}(-\theta^\ast)
\]
for the boundary values obtained from real rapidity.  Crossing for an
self-conjugate scalar requires \(S_{11}(\theta)=S_{11}(\ii\pi-\theta)\).  Since
\(\sinh(\ii\pi-\theta)=\sinh\theta\), the displayed amplitude satisfies this
identity.

The denominator vanishes when
\[
  \sinh\theta=\ii\sin(\pi\xi).
\]
Inside the physical strip \(0<\operatorname{Im}\theta<\pi\), this equation has
two solutions,
\[
  \theta=\ii\pi\xi,
  \qquad
  \theta=\ii\pi(1-\xi).
\]
They are exchanged by crossing.  Only the first one is the direct
\(B_1B_1\to B_2\) pole when \(2\xi<1\); the second is its crossed-channel
image and must not be counted as a separate direct-channel particle.
For two incoming particles of mass \(m_1\) and rapidities
\(\theta/2\) and \(-\theta/2\), the squared center-of-mass energy is
\[
  s(\theta)=
  \bigl(p(\theta/2)+p(-\theta/2)\bigr)^2
  =
  4m_1^2\cosh^2(\theta/2).
\]
At \(\theta=\ii\pi\xi\), this is
\[
  s_\ast
  =
  4m_1^2\cos^2\!\left(\frac{\pi\xi}{2}\right).
\]
Therefore the pole has the mass required for the second breather:
\[
  m_2=2m_1\cos\!\left(\frac{\pi\xi}{2}\right).
\]
Using the soliton-breather mass formula
\[
  m_n=2M_s\sin\!\left(\frac{\pi n\xi}{2}\right)
\]
for \(n=1,2\), the same relation becomes
\[
  2m_1\cos\!\left(\frac{\pi\xi}{2}\right)
  =
  4M_s\sin\!\left(\frac{\pi\xi}{2}\right)
      \cos\!\left(\frac{\pi\xi}{2}\right)
  =
  2M_s\sin(\pi\xi)
  =
  m_2 .
\]
Thus the elementary analytic pole of
\eqref{eq:sine-gordon-b1b1-smatrix} encodes the fusion process
\[
  B_1+B_1\longrightarrow B_2
\]
when \(B_2\) lies below the two-soliton threshold as a stable breather.  The
condition is \(2\xi<1\).  At the endpoint \(2\xi=1\) the candidate sits at
threshold, and for \(2\xi>1\) this pole does not define an additional stable
breather in the Hilbert-space particle spectrum.

\section{Neutral Breather Blocks and Minimal Fusion}

It is useful to isolate the scalar factor
\begin{equation}
  F_a(\theta)
  =
  \frac{\sinh\theta+\ii\sin(\pi a)}
       {\sinh\theta-\ii\sin(\pi a)} ,
  \qquad a\in\mathbb R .
  \label{eq:sine-gordon-neutral-block}
\end{equation}
Then \(S_{11}(\theta)=F_\xi(\theta)\).  This is not merely notation: the
block records both the direct pole and its crossed image in a way that makes
the residue sign visible.

\begin{proposition}[Neutral scalar block]
\label{prop:sine-gordon-neutral-block}
For real \(0<a<1\),
\[
  F_a(\theta)F_a(-\theta)=1,
  \qquad
  F_a(\ii\pi-\theta)=F_a(\theta).
\]
The physical strip contains the two simple poles
\[
  \theta=\ii\pi a,
  \qquad
  \theta=\ii\pi(1-a).
\]
The corresponding residues obey
\[
  -\ii\,\operatorname*{Res}_{\theta=\ii\pi a}F_a(\theta)
  =
  2\tan(\pi a),
  \qquad
  -\ii\,\operatorname*{Res}_{\theta=\ii\pi(1-a)}F_a(\theta)
  =
  -2\tan(\pi a).
\]
Thus, for \(0<a<1/2\), the pole at \(\theta=\ii\pi a\) has the positive
direct-channel residue sign, while the pole at \(\theta=\ii\pi(1-a)\) is its
crossed-channel partner.
\end{proposition}

\begin{proof}
Unitarity follows by replacing \(\sinh\theta\) by \(-\sinh\theta\), and
crossing follows from \(\sinh(\ii\pi-\theta)=\sinh\theta\).  The denominator
vanishes at \(\theta=\ii\pi a\) and \(\theta=\ii\pi(1-a)\).  At
\(\theta_0=\ii\pi a\), the numerator is \(2\ii\sin(\pi a)\), while the
derivative of the denominator is \(\cosh\theta_0=\cos(\pi a)\).  Therefore
\[
  \operatorname*{Res}_{\theta=\ii\pi a}F_a(\theta)
  =
  2\ii\tan(\pi a).
\]
At \(\theta_1=\ii\pi(1-a)\), the numerator is again
\(2\ii\sin(\pi a)\), but the derivative is
\(\cos(\pi(1-a))=-\cos(\pi a)\), giving the second residue formula.
\end{proof}

The minimal diagonal scattering amplitude of breathers \(B_m\) and \(B_n\),
when both labels and every displayed fused label lie in the stable range
\(k\xi<1\), is
\begin{equation}
  S_{mn}(\theta)
  =
  F_{\frac{(m+n)\xi}{2}}(\theta)\,
  F_{\frac{|m-n|\xi}{2}}(\theta)\,
  \prod_{\ell=1}^{\min(m,n)-1}
  F_{\frac{(|m-n|+2\ell)\xi}{2}}(\theta)^2 ,
  \label{eq:sine-gordon-breather-breather-smatrix}
\end{equation}
with the convention \(F_0(\theta)=1\).  The formula is the minimal solution
of the diagonal bootstrap equations with the breather spectrum fixed above:
each \(F_a\) supplies one direct pole and its crossing partner, and the
double factors are exactly the poles encountered twice when the two fused
constituent rapidity strings scatter through each other.  More explicitly,
start from \(S_{11}=F_\xi\).  The residue identity of Proposition
\ref{prop:fusion-identity-residue-derivation} constructs \(B_{r+s}\) from
the direct pole of \(B_rB_s\); because the neutral breather amplitudes are
diagonal scalar functions, the fused scattering with a third breather is the
ordinary product of the constituent scalar factors evaluated at the shifted
rapidities.  Iterating this statement produces the two endpoint factors
\((m+n)\xi/2\) and \(|m-n|\xi/2\), while every intermediate constituent
difference \((|m-n|+2\ell)\xi/2\) occurs once from each side of the fused
rapidity string and therefore appears squared.  The word minimal means that
no further CDD factor, analytic and pole-free in the physical strip, has been
multiplied into the result.

\begin{proposition}[Direct fusion pole in the breather amplitudes]
\label{prop:sine-gordon-breather-direct-fusion}
Assume \(m,n\ge1\) and \((m+n)\xi<1\), so that \(B_{m+n}\) is a stable
breather.  The amplitude \(S_{mn}(\theta)\) in
\eqref{eq:sine-gordon-breather-breather-smatrix} has a direct-channel pole at
\[
  \theta=\ii u_{mn}^{m+n},
  \qquad
  u_{mn}^{m+n}=\frac{\pi(m+n)\xi}{2},
\]
and the bound-state kinematics at this pole gives the mass \(m_{m+n}\).
\end{proposition}

\begin{proof}
The pole is supplied by the first factor
\(F_{(m+n)\xi/2}(\theta)\).  Since \((m+n)\xi<1\), its parameter
\((m+n)\xi/2\) lies in \((0,1/2)\), so Proposition
\ref{prop:sine-gordon-neutral-block} gives the positive direct-channel
residue sign at the displayed rapidity.  Let
\[
  \alpha=\frac{\pi\xi}{2},
  \qquad
  m_k=2M_s\sin(k\alpha).
\]
The pole angle is \(u=(m+n)\alpha\).  The rapidity kinematics gives
\[
  m_m^2+m_n^2+2m_m m_n\cos u
  =
  4M_s^2\left[
  \sin^2(m\alpha)+\sin^2(n\alpha)
  +2\sin(m\alpha)\sin(n\alpha)\cos((m+n)\alpha)
  \right].
\]
Using
\[
  \sin((m+n)\alpha)
  =
  \sin(m\alpha)\cos(n\alpha)+\cos(m\alpha)\sin(n\alpha),
\]
one obtains the elementary identity
\[
  \sin^2((m+n)\alpha)
  =
  \sin^2(m\alpha)+\sin^2(n\alpha)
  +2\sin(m\alpha)\sin(n\alpha)\cos((m+n)\alpha).
\]
Thus the bracket in the preceding display is \(\sin^2((m+n)\alpha)\), and
the invariant mass at the pole is
\[
  2M_s\sin((m+n)\alpha)=m_{m+n},
\]
which proves the claim.
\end{proof}

\section{Massive Thirring Equivalence}

The massive Thirring model has Dirac field \(\psi\) and action
\[
  S_{\mathrm{MT}}
  =
  \int d^2x\,
  \left[
  \ii\bar\psi\gamma^\mu\partial_\mu\psi
  -m\bar\psi\psi
  -\frac{g_{\mathrm{T}}}{2}
  (\bar\psi\gamma^\mu\psi)(\bar\psi\gamma_\mu\psi)
  \right].
\]
In Coleman's normalization, the sine-Gordon and massive Thirring couplings
are related by
\[
  \frac{4\pi}{\beta^2}=1+\frac{g_{\mathrm{T}}}{\pi}.
\]
The equivalence identifies the Thirring fermion number with the sine-Gordon
topological charge:
\[
  \bar\psi\gamma^\mu\psi
  =
  \frac{\beta}{2\pi}\epsilon^{\mu\nu}\partial_\nu\varphi .
\]
At the free-fermion point \(\beta^2=4\pi\), hence \(g_{\mathrm{T}}=0\), the
soliton and antisoliton form the massive Dirac particle and antiparticle.
Away from this point the equivalence is a statement about renormalized local
operator algebras and scattering data in the declared convention.

\section{Affine Toda Data}

Let \(\mathfrak g\) be a simple Lie algebra of rank \(r\), with simple roots
\(\alpha_i\) and highest root
\[
  \theta=\sum_{i=1}^r n_i\alpha_i .
\]
Set \(\alpha_0=-\theta\) and \(n_0=1\).  The affine Toda action for a real
Cartan-valued field \(\phi\in\mathfrak h_\mathbb R\) is
\[
  S_{\mathrm{AT}}
  =
  \int d^2x\,
  \left[
  -\frac12(\partial_\mu\phi,\partial^\mu\phi)
  -\frac{m^2}{\beta^2}\sum_{i=0}^r n_i
  \left(e^{\beta(\alpha_i,\phi)}-1\right)
  \right].
\]
The classical equations admit a Lax representation with spectral parameter.
Consequently the classical theory has commuting local conserved charges.  In
the quantum theory, the factorized S-matrix is organized by the root system,
particle species are labelled by nodes of the Dynkin diagram in the
simply-laced cases, and the pole structure encodes fusing rules compatible
with the affine root data.

\section{Mass Matrix In Affine Toda Theory}

Expanding the potential at \(\phi=0\) gives
\[
  V(\phi)
  =
  \frac{m^2}{2}
  \sum_{i=0}^r n_i(\alpha_i,\phi)^2+O(\phi^3).
\]
Thus the classical mass matrix is
\[
  M_{\mathrm{cl}}^2
  =
  m^2\sum_{i=0}^r n_i\,\alpha_i\otimes\alpha_i .
\]
For simply-laced algebras its eigenvectors are related to the Perron-Frobenius
eigenvector of the Cartan matrix, and the mass ratios match the Coxeter data
entering the exact scattering blocks.  A quantum statement about mass ratios
requires fixing the coupling normalization and the renormalized mass
coordinate.

\section{Model-Comparison Ledger}

The exact-solution claim for each model has three independent tests:
\[
\begin{array}{ll}
\text{UV test} & \text{operator dimensions and perturbing-coordinate
normalization},\\
\text{IR test} & \text{particle masses, charges, and factorized S-matrix
poles},\\
\text{finite-size test} & \text{TBA ground-state energy and local operator
finite-volume data}.
\end{array}
\]
Agreement of one row does not determine the other rows.  The model becomes a
QFT example for this monograph only after all rows are expressed in the same
regularization and source convention.

\begin{openproblem}[Constructive integrable model comparison]
\label{op:constructive-integrable-model-comparison}
Construct sine-Gordon and affine Toda theories as continuum local QFTs with
operators, soliton sectors, factorized scattering, TBA finite-size energies,
and form factors derived from one regulated framework rather than matched
only through separate exact-solution data.
\end{openproblem}
